\begin{lemma}[Integration by parts for a $C^1$ composition against a test function]
Fix $n \in \mathbb{N}_0$ and write $\lambda$ for Lebesgue measure on
$\mathbb{R}^n$.  Let $q, q' : \mathbb{R} \to \mathbb{R}$ be such that $q$ is
differentiable at every $t \in \mathbb{R}$ with derivative $q'(t)$, and such
that $q'$ is continuous, so $q \in C^1(\mathbb{R})$.  Fix a coordinate index
$j \in \{1,\ldots,n\}$, and let $w : \mathbb{R}^n \to \mathbb{R}$ be
differentiable at every point of $\mathbb{R}^n$ and such that its $j$th partial
derivative $x \mapsto \partial_j w(x)$ is continuous on $\mathbb{R}^n$.  Let
$\varphi : \mathbb{R}^n \to \mathbb{R}$ be smooth with compact support.  Then
\[
  \int_{\mathbb{R}^n} q(w(x))\, \partial_j \varphi(x)\, d\lambda(x)
    = -\int_{\mathbb{R}^n}
        \bigl(q'(w(x))\, \partial_j w(x)\bigr)\, \varphi(x)\, d\lambda(x),
\]
where $\partial_j$ denotes the Fr\'echet derivative evaluated at the $j$-th
standard basis vector $e_j$.  Both integrals converge, and the identity is one
of real numbers.

\emph{The hypothesis on $w$ is $C^1$ in the $j$th direction, and no more.}
Everywhere differentiability plus continuity of the single partial
$\partial_j w$ is exactly what the argument below consumes; neither higher
smoothness of $w$ nor continuity of the full Fr\'echet derivative $Dw$ is used,
and neither is assumed.  That is why the index $j$ is quantified ahead of the
hypothesis rather than after it: for each separate direction the identity has
its own, separate, minimal hypothesis.  A smooth $w$ satisfies both conditions
for every $j$, so the smooth case is contained in this one.

Every integral in this statement is taken against Lebesgue measure.  No Gaussian
weight appears: although this lemma is used inside the tablet's Gaussian
development, the statement and its proof are purely classical and involve
$\gamma_n$ nowhere.  Correspondingly it cites no other tablet node at all: it
rests only on the shared background calculus, which is what makes it a
self-contained classical identity worth isolating.
\end{lemma}
\begin{proof}
Abbreviate $f = q \circ w$, so $f(x) = q(w(x))$.

Step 1 (the chain rule for $f$).  By hypothesis $w$ is differentiable at every
$x$, with Fr\'echet derivative $Dw(x)$, and $q$ is differentiable at every real
number, so the composite is differentiable at every $x$ with
\[
  Df(x) = q'(w(x))\, Dw(x)
\]
as continuous linear maps $\mathbb{R}^n \to \mathbb{R}$, the scalar
$q'(w(x))$ multiplying the linear map $Dw(x)$.  Evaluating at $e_j$ gives the
pointwise identity
\[
  \partial_j f(x) = q'(w(x))\, \partial_j w(x)
  \qquad \text{for every } x \in \mathbb{R}^n . \tag{C}
\]

Step 2 (continuity of the four integrands' factors).  $q$ is continuous, being
differentiable everywhere, and $w$ is continuous, being differentiable
everywhere, so $f = q \circ w$ is continuous as a composite of continuous maps.
The map $x \mapsto \partial_j w(x)$ is continuous by hypothesis, and
$x \mapsto \partial_j \varphi(x)$ is continuous because $\varphi$ is smooth, so
its derivative is continuous and evaluation at $e_j$ is continuous.  By (C),
$x \mapsto \partial_j f(x)$ is then continuous as well, being the product of the
continuous $q' \circ w$ with the continuous $\partial_j w$.  Finally
$\partial_j \varphi$ has compact support, the support of a derivative being
contained in the support of the function.

Step 3 (integrability of the three products).  Integration by parts requires the
three products
\[
  \partial_j f \cdot \varphi, \qquad f \cdot \partial_j \varphi, \qquad
  f \cdot \varphi
\]
to be $\lambda$-integrable.  Each is a continuous function with compact support:
in the first and third the factor $\varphi$ has compact support, and in the
second the factor $\partial_j \varphi$ does, while the other factor is
continuous by Step 2.  A continuous function with compact support is
$\lambda$-integrable, so all three integrals converge.

Step 4 (integration by parts).  Both $f$ and $\varphi$ are differentiable at
every point, in particular at every point of the closed support of the other, so
the classical integration-by-parts identity in the direction $e_j$ applies to
the pair $(f, \varphi)$ with the integrability of Step 3 as its hypothesis, and
yields
\[
  \int_{\mathbb{R}^n} f(x)\, \partial_j \varphi(x)\, d\lambda(x)
    = -\int_{\mathbb{R}^n} \partial_j f(x)\, \varphi(x)\, d\lambda(x) .
\]
There is no boundary term because $\varphi$ is compactly supported.

Step 5 (substituting the chain rule).  By (C) the integrands on the right agree
at every point:
$\partial_j f(x)\, \varphi(x)
   = \bigl(q'(w(x))\, \partial_j w(x)\bigr)\varphi(x)$ for every $x$, hence in
particular $\lambda$-almost everywhere, so the two integrals are equal.
Substituting into Step 4 and recalling $f = q \circ w$ gives
\[
  \int_{\mathbb{R}^n} q(w(x))\, \partial_j \varphi(x)\, d\lambda(x)
    = -\int_{\mathbb{R}^n}
        \bigl(q'(w(x))\, \partial_j w(x)\bigr) \varphi(x)\, d\lambda(x),
\]
which is the assertion.
\end{proof}
